\begin{problem}{}{}{Real Mathematical Analysis, Pugh, Ex.9 p.198}
Assume that $f : \mathbb{R} \to \mathbb{R}$ is differentiable.
\begin{enumerate}
    \item[(a)] If there is an $L < 1$ such that for each $x \in \mathbb{R}$ we have $f'(x) < L$, prove that there exists a unique point $x$ such that $f(x) = x$. [$x$ is a fixed point for $f$.]
    \item[(b)] Show by example that (a) fails if $L = 1$.
\end{enumerate}
\end{problem}

\begin{solution}
    \begin{itemize}
        \item[(a)]
              We have $f'(x) < L$, $\forall x \in \mathbb{R}$.
              This implies that $\frac{f(a) - f(b)}{a - b} < L$, $\forall a,b \in \mathbb{R}, a < b$.
              Otherwise, $\exists a, b \in \mathbb{R}, a < b : \frac{f(a) - f(b)}{a - b} \ge L$.
              By MVT, $\exists \theta \in (a, b) : f'(\theta) = \frac{f(a) - f(b)}{a - b} \ge L$,
              which is a contradiction.

              Define $g(x) = f(x) - x$. We have:
              \begin{align*}
                           & \frac{f(a) - f(b)}{a - b} < L                                         \\
                  \implies & f(a) - f(b) > L(a - b)                                                \\
                  \implies & g(a) - g(b) > (L - 1)(a - b), \quad \forall a,b \in \mathbb{R}, a < b
              \end{align*}

              Plug in $b = 0 \implies g(a) > g(0) + (L - 1)a$, $\forall a < 0$. \\
              Since $L < 1$ and $a < 0$, we have $(L - 1)a > 0$ and this can be made
              arbitrarily large by letting $a \to -\infty$.
              This implies that $\exists a_1 < 0$ such that $g(a_1) > 0$.

              Plug in $a = 0 \implies g(b) < (L - 1)b + g(0)$, $\forall b > 0$. \\
              Similarly, we can make $g(b)$ be arbitrarily small by letting $b \to \infty$.
              This implies $\exists b_1 > 0 : g(b_1) < 0$.

              Since $g$ is differentiable on $\mathbb{R}$, $g$ is also continuous and
              therefore has the Intermediate Value Property.
              There exists $\theta \in (a_1, b_1)$ such that $g(\theta) = 0$ or $f(\theta) = \theta$.

              Now suppose $\exists \theta_1, \theta_2 \in \mathbb{R}$, $\theta_1 \neq \theta_2$
              such that $f(\theta_1) = \theta_1$ and $f(\theta_2) = \theta_2$.
              It follows that:
              \[
                  1 = \frac{f(\theta_1) - f(\theta_2)}{\theta_1 - \theta_2} < L
              \]
              which is a contradiction. Thus, the fixed point is unique.
        \item[(b)]
              Let $L = 1$. We know that $g(x)$ is decreasing on $\mathbb{R}$ since
              $g'(x) = f'(x) - 1 < 0$, $\forall x \in \mathbb{R}$.
              We want to find an example that fails to have a root, meaning $g$ needs
              to be decreasing but never cross $y = 0$.

              Let $g(x) = e^{-x}$. We then have:
              \begin{align*}
                  f(x)           & = e^{-x} + x                                      \\
                  \implies f'(x) & = 1 - e^{-x} < 1, \quad \forall x \in \mathbb{R},
              \end{align*}
              which satisfies our desired condition.

              Now, checking for a fixed point:
              \begin{align*}
                           & f(x) = x       \\
                  \implies & e^{-x} + x = x \\
                  \implies & e^{-x} = 0
              \end{align*}

              But $e^{-x} > 0$, $\forall x \in \mathbb{R}$, so there is no fixed point.
    \end{itemize}
\end{solution}